```latex
\section{Результаты}

В результате выполнения лабораторной работы была реализована программа,
состоящая из двух отдельных процессов: родительского и дочернего.

Родительский процесс создаёт неименованный канал и дочерний процесс.
Читающий конец канала перенаправляется на стандартный поток ввода
дочернего процесса. После этого родительский процесс передаёт через канал
введённые пользователем строки, а дочерний процесс вычисляет сумму чисел
в каждой строке и записывает результат в файл.

При выполнении программы в консоль выводятся PID родительского и
дочернего процессов, что позволяет убедиться в создании отдельного
дочернего процесса.

Пример входных данных:

\begin{lstlisting}
result.txt
2.3 3.3 5 6.6 102.432
3.4 5.4 7.6
4.5
\end{lstlisting}

После завершения ввода в файл \texttt{result.txt} записываются суммы
соответствующих строк:

\begin{lstlisting}
119.632004
16.400000
4.500000
\end{lstlisting}

Работа программы была дополнительно исследована с помощью утилиты
\texttt{strace}. В трассировке были обнаружены системные вызовы,
соответствующие основным этапам работы программы:

\begin{itemize}
    \item \texttt{pipe2()} --- создание неименованного канала;
    \item \texttt{clone()} --- создание дочернего процесса;
    \item \texttt{dup2()} --- перенаправление читающего конца канала
    на стандартный ввод дочернего процесса;
    \item \texttt{execve()} --- запуск программы \texttt{child};
    \item \texttt{write()} и \texttt{read()} --- передача данных
    через канал;
    \item \texttt{openat()} --- открытие выходного файла;
    \item \texttt{wait4()} --- ожидание завершения дочернего процесса.
\end{itemize}

После завершения пользовательского ввода родительский процесс закрывает
пишущий конец канала. Дочерний процесс получает EOF, записывает
накопленные результаты в файл, закрывает файл и завершается. Родительский
процесс получает статус завершения дочернего процесса и также завершает
работу.